% Options for packages loaded elsewhere
\PassOptionsToPackage{unicode}{hyperref}
\PassOptionsToPackage{hyphens}{url}
\documentclass[
]{book}
\usepackage{xcolor}
\usepackage{amsmath,amssymb}
\setcounter{secnumdepth}{5}
\usepackage{iftex}
\ifPDFTeX
  \usepackage[T1]{fontenc}
  \usepackage[utf8]{inputenc}
  \usepackage{textcomp} % provide euro and other symbols
\else % if luatex or xetex
  \usepackage{unicode-math} % this also loads fontspec
  \defaultfontfeatures{Scale=MatchLowercase}
  \defaultfontfeatures[\rmfamily]{Ligatures=TeX,Scale=1}
\fi
\usepackage{lmodern}
\ifPDFTeX\else
  % xetex/luatex font selection
\fi
% Use upquote if available, for straight quotes in verbatim environments
\IfFileExists{upquote.sty}{\usepackage{upquote}}{}
\IfFileExists{microtype.sty}{% use microtype if available
  \usepackage[]{microtype}
  \UseMicrotypeSet[protrusion]{basicmath} % disable protrusion for tt fonts
}{}
\makeatletter
\@ifundefined{KOMAClassName}{% if non-KOMA class
  \IfFileExists{parskip.sty}{%
    \usepackage{parskip}
  }{% else
    \setlength{\parindent}{0pt}
    \setlength{\parskip}{6pt plus 2pt minus 1pt}}
}{% if KOMA class
  \KOMAoptions{parskip=half}}
\makeatother
\usepackage{longtable,booktabs,array}
\usepackage{calc} % for calculating minipage widths
% Correct order of tables after \paragraph or \subparagraph
\usepackage{etoolbox}
\makeatletter
\patchcmd\longtable{\par}{\if@noskipsec\mbox{}\fi\par}{}{}
\makeatother
% Allow footnotes in longtable head/foot
\IfFileExists{footnotehyper.sty}{\usepackage{footnotehyper}}{\usepackage{footnote}}
\makesavenoteenv{longtable}
\usepackage{graphicx}
\makeatletter
\newsavebox\pandoc@box
\newcommand*\pandocbounded[1]{% scales image to fit in text height/width
  \sbox\pandoc@box{#1}%
  \Gscale@div\@tempa{\textheight}{\dimexpr\ht\pandoc@box+\dp\pandoc@box\relax}%
  \Gscale@div\@tempb{\linewidth}{\wd\pandoc@box}%
  \ifdim\@tempb\p@<\@tempa\p@\let\@tempa\@tempb\fi% select the smaller of both
  \ifdim\@tempa\p@<\p@\scalebox{\@tempa}{\usebox\pandoc@box}%
  \else\usebox{\pandoc@box}%
  \fi%
}
% Set default figure placement to htbp
\def\fps@figure{htbp}
\makeatother
\setlength{\emergencystretch}{3em} % prevent overfull lines
\providecommand{\tightlist}{%
  \setlength{\itemsep}{0pt}\setlength{\parskip}{0pt}}
\usepackage[]{natbib}
\bibliographystyle{plainnat}
\usepackage{booktabs}
\usepackage{bookmark}
\IfFileExists{xurl.sty}{\usepackage{xurl}}{} % add URL line breaks if available
\urlstyle{same}
\hypersetup{
  pdftitle={Psykiatrian tärpit},
  pdfauthor={Remo Kapanen},
  hidelinks,
  pdfcreator={LaTeX via pandoc}}

\title{Psykiatrian tärpit}
\author{Remo Kapanen}
\date{2026-03-01}

\begin{document}
\maketitle

{
\setcounter{tocdepth}{1}
\tableofcontents
}
\chapter{About}\label{about}

Tentit koostuvat avonaisista kysymyksistä: niin hieman suljetuimmista kysymyksistä kuin laajoista potilastapauksista.

Suuressa osassa tenteistä ei ole mallivastauksia, joten olen näiden kohdalla yrittänyt vastata omalla tavallani kysymyksiin. Joissain taas löytyy mallivastaukset, mutta ne ovat usein vaikealukuisia ja näiden kohdalla olen yrittänyt muokata vastauksia siistimmiksi.

\chapter{2026}\label{section}

\section{Blokki 1}\label{blokki-1}

\subsection{Unettomuus TK-vastaanotolla (5 pistettä)}\label{unettomuus-tk-vastaanotolla-5-pistettuxe4}

52-vuotias mies saapuu terveyskeskukseen vastaanotollesi unettomuuden vuoksi. Hänellä on pitkäaikaissairautena verenpainetauti, obesiteetti ja lisäksi hän kärsii kroonisista selkäkivuista. Potilas kertoo että hänellä on puolen vuoden ajan ollut vaikea nukahtaa ja öisin hän heräilee eikä saa enää unta. Puoliso on sanonut että potilas on ollut jo pidempään vain ''varjo entisestään''. Töissä on vaikea jaksaa olla ja keskittyä sekä töiden jälkeen on kovin uupunut. Harrastukset on jääneet hiljattain pois. Vuorokausirytmi on normaali ja potilas ei käytä alkoholia. Potilas kuvailee iltoja niin että vaikea rauhoittua nukkumaan, pitää lähteä ''takaisin tuvan puolelle''. Kertoo ohimennen että vuosi sitten tyttärellä todettiin leukemia, hoidot on olleet käynnissä ja tytär voi hyvin. Potilas on saanut työterveyslääkäriltä aikaisemmin doksepiini-reseptin unettomuuteen mutta epäröi sitä aloittaa.

\begin{itemize}
\tightlist
\item
  \begin{enumerate}
  \def\labelenumi{\alph{enumi})}
  \tightlist
  \item
    Mitä diagnostisia vaihtoehtoja lähtisit näillä tiedoin lisäselvittämään vastaanoton aikana unettomuushäiriön lisäksi? Mainitse vähintään neljä.
  \end{enumerate}
\item
  \begin{enumerate}
  \def\labelenumi{\alph{enumi})}
  \setcounter{enumi}{1}
  \tightlist
  \item
    Potilas kysyy ''onko doksepiinia turvallista käyttää ja mitä se tekee''. Mitä vastaat hänelle?
  \end{enumerate}
\item
  \begin{enumerate}
  \def\labelenumi{\alph{enumi})}
  \setcounter{enumi}{2}
  \tightlist
  \item
    Potilas on menossa ensikäynnille unihoitajalle jonka luona olisi tarkoitus aloittaa CBT-I hoitojakso unettomuuden hoitoon. Hän kysyy sinulta ''mistä siinä hoidossa oikein on kyse''. Mitä vastaat hänelle?
  \end{enumerate}
\end{itemize}

a

Ei-elimellinen unettomuus (F51.0) on diagnoosi, jota tulisi käyttää, jos potilaan vallitseva oire on unen epätyydyttävä laatu tai määrä eikä univaikeus selity muulla sairaudella. Unettomuushäiriön muut diagnostiset kriteerit kyllä täyttyvät potilaalla. Diagnostisia kriteerejä siis ovat:

A: Vallitseva oire on tyytymättömyys unen määrään tai laatuun ja tämän lisäksi yksi tai useampi seuraavista oireista:

\begin{enumerate}
\def\labelenumi{\arabic{enumi}.}
\tightlist
\item
  Nukahtamisvaikeus (lapsilla tämä voi ilmetä vaikeutena nukahtaa ilman huoltajan apua)

  Unessa pysymisen vaikeus, joka ilmenee toistuvina heräämisinä tai vaikeutena nukahtaa yöllä uudelleen heräämisen jälkeen

  Herääminen aamulla liian aikaisin kykenemättä enää nukahtamaan uudelleen
\end{enumerate}

B: Unihäiriö aiheuttaa kliinisesti merkittävää kärsimystä tai haittaa jollakin tärkeällä toimintakyvyn osa-alueella (esim. sosiaalinen, ammatillinen, kasvatuksellinen, akateeminen toimintakyky tai käyttäytyminen)

C: Univaikeus esiintyy vähintään kolmena yönä viikossa vähintään kuukauden ajan*

D: Häiriötä ei aiheuta elimellinen, esimerkiksi neurologinen, syy tai sisätautiongelma, psyykkisiin toimintoihin vaikuttava lääkitys tai muu lääkitys

Diagnostiikkaan siis kuuluu muiden syiden poissulku ja tässä huolellinen anamneesi on tärkein työkalu. Tehtävänannosta lukien potilaalla on verenpainetauti, obesiteetti ja lisäksi hän kärsii kroonisista selkäkivuista. Hän myös kuvailee väsymystä päivällä: ``töiden jälkeen on kovin uupunut''.

\begin{center}\rule{0.5\linewidth}{0.5pt}\end{center}

\textbf{Erotusdiagnostinen vaihtoehto 1: Kipu}

Potilashistoriasta lukien hän kärsii kroonisista selkäkivuista. Kipu on todella yleinen unettomuuden syy -- selkäkivut voivat estää rentoutumisen ja aiheuttaa heräilyä kivun vuoksi. Kivun luonteesta ja sen yhteydestä uniongelmiin tulisi kysyä.

\textbf{Erotusdiagnostinen vaihtoehto 2: Masennus}

Potilaan puoliso kuvailee häntä ``varjoksi entisestään'', harrastukset ovat jääneet ja hän kärsii uupumuksesta sekä keskittymisvaikeuksista. Nämä ovat tyypillisiä masennuksen oireita ja masennuksen klassiselle muodolle on tyypillistä uniongelmat.

Masennusoireista tulee kysellä enemmän ja selvittää, sopisiko oirekuva ihan masennustilaksi.

\textbf{Erotusdiagnostinen vaihtoehto 3: Uniapnea}

Uniapnea voi ilmetä unettomuutena (potilas herää hengityshäiriötapahtumiin).

STOP-BANG-pisteikköä käyttäen potilaalle tulee tarpeeksi pisteitä, että sen mahdollisuus tulisi huomioida. Potilas on yli 50-vuotias (1p) ja mies (1p), hän on ylipainoinen (mahdollisesti 1p, BMI:tä ei ole kerrottu), hänellä on korkea verenpaine (1p) ja kokee itsensä väsyneeksi päiväaikaan (1p) -\textgreater{} yhteensä 4-5 pistettä pelkästään tehtävänannon perusteella ja uniapnean todennäköisyys on suuri, jos STOP-BANG on 5-8 pistettä.

Potilaalta tulisi siis kysellä, että kuorsaako äänekkäästi ja kysyä onko hänen vaimonsa koskaan huomauttanut hengityskatkoksista/huomannut näitä itse. Mahdollisesti potilaalle voisi algoritmisesti teettää yöpolygrafian, vaikka hengityskatkoksia ei oltaisi huomattu.

\textbf{Erotusdiagnostinen vaihtoehto 4: Levottomat jalat -oireyhtymä (RLS)}

Potilas kuvailee iltoja niin että vaikea rauhoittua nukkumaan ja pitää lähteä ''takaisin tuvan puolelle''.

Levottomien jalkojen oireyhtymällä tarkoitetaan alaraajoissa levon aikana (illalla tai yöllä) esiintyvää, vaikeasti kuvattavaa epämiellyttävää tuntemusta, joka helpottuu alaraajojen liikuttelulla. Potilaan oirekuva siis sinänsä sopii tähän. Tuntemuksista nukkumaan menon yhteydessä tulee kysyä tarkemmin ja arvioida, sopivatko oireet RLS:ään.

Potilaalla on myös diabetes, joten jos oirekuva sopisi RLS:ään, tulee kliinisellä tutkimuksella (turvotukset, suonikohjut, staasi-ihottuma, alaraajojen valtimoverenkierto, kosketus-, terävä- ja värinätunnot, lihaskato, lihasjänteys) sulkea pois muut sairaudet, kuten polyneuropatia, valtimo- ja laskimoinsuffisienssi.

\textbf{Erotusdiagnostinen vaihtoehto 5: Ahdistuneisuushäiriö}

Vaikea rauhoittua nukahtamaan ja keskittymishäiriöitä töissä. Periaatteessa nämä sopivat moniin ahdistuneisuushäiriöihin, erityisesti yleistyneeseen ahdistuneisuushäiriöön (GAD). Kannattaa kysellä yleisemmin potilaan ahdistusoireista.

\textbf{Erotusdiagnostinen vaihtoehto 6: Lääkitys/elämäntavat/unihygienia}

Potilaalla on sairauksia ja mahdollisesti lääkityksiä niihin (tehtävänannossa ei mainita lääkelistaa). Potilas saattaa syödä jotain unettomuutta aiheuttavaa lääkettä, kannattaa tarkastaa lääkelista.

Tehtävänannossa mainitaan, että vuorokausirytmi on normaali ja potilas ei käytä alkoholia. Tämä ei kuitenkaan poissulje sitä, että potilas esimerkiksi joisi todella paljon kahvia myöhään päivälläkin, joten esim. siitä kannattaa kysyä. Samoin tulee kysyä muutenkin yleisesti unihygieniasta ja esimerkiksi selaako stimuloivaa materiaalia kännykällä sängyssä.

\subsection{Keskittymisvaikeuksia ja opintosuoritusten heikkenemistä opiskelijalla (5 pistettä)}\label{keskittymisvaikeuksia-ja-opintosuoritusten-heikkenemistuxe4-opiskelijalla-5-pistettuxe4}

22‑vuotias ammattiopiston opiskelija Oskari tulee vastaanotollesi viime kuukausina lisääntyneiden keskittymisvaikeuksien ja opintosuoritusten heikkenemisen vuoksi. Hän kertoo, että opinnoissa ``ajatus ei oikein pysy kasassa'' ja keskustelutilanteissa ``sanoja on vaikea löytää''. Ollut yhteydessä opiskelijahuollon terveydenhoitajaan, joka ohjannut Oskarin vastaanotollesi perusterveydenhuollon yksikköön.

Erityisesti viimeisen kahden kuukauden aikana Oskari on alkanut huomata seuraavaa:

\begin{itemize}
\tightlist
\item
  Ajatusten omituisia katkoksia, jolloin ``mieli menee tyhjäksi''.
\item
  Epämääräistä epävarmuutta siitä, onko joku maininnut hänen nimensä julkisissa paikoissa, mutta hän ei ole varma, onko kuullut oikein.
\item
  Hän kertoo myös, että ``oma ajattelu ei ole entisenlaista'', ja hänen on vaikea luottaa siihen, että omat havainnot vastaavat todellisuutta.
\item
  Ystävät ovat huomauttaneet, että hän vaikuttaa vetäytyneeltä ja poissaolevalta.
\item
  Mieliala on ajoittain matalampi, mutta ei osaa tätä tarkemmin kuvata. Oskari kertoo, ettei hänellä ole aikaisempia lääkärikäyntejä. Alkoholia käyttää satunnaisesti.
\item
  Opiskelijajuhlissa kokeillut kerran kannabista, josta tuli ``huono olo''.
\end{itemize}

Oskari on huolissaan toimintakykynsä laskusta, mutta ei itse osaa selittää, mikä on muuttunut. \textbf{Miten lähdet tutkimaan ja selvittämään Oskarin tilannetta? Pohdi diagnostisia mahdollisuuksia. Jatkotoimenpiteesi selvitystesi perusteella?}

\subsection{(5 pistettä)}\label{pistettuxe4}

Vastaanotollesi terveyskeskukseen tulee 56-vuotias nainen, jolle on kertomansa mukaan ilmaantunut viimeisten kuukausien aikana voimakasta väsymystä ja alentunutta toimintakykyä. Aiempien sairauskertomustekstien perusteella hän on osa-aikainen perhepäivähoitaja,kahden nuoren aikuisen vanhempi ja avioliitossa.
Aluksi hän kertoo, ettei muista mitään eikä oikein saa tehdyksi mitään. Hän kertoo kokeneensa ajoittain mielialansa oikein surulliseksikin, mutta välillä olo on ollut täysin turta. Kertoo, että on vielä loppusyksystä touhunnut kaikenlaista harrastustoimintaa, mutta nyt kaikki on jäänyt. Töistä on suoriutunut ``joten kuten'', mutta työpäivien jälkeen kertomansa mukaan on vain maannut kotona ja ruokahalu on ollut huono. Kertoo tunteneensa itsensä jatkuvasti epäonnistuneeksi.
Potilas on ulkoiselta olemukseltaan siisti ja keskustelussa vastavuoroisesti toimiva. Hänen kerrontansa on loogista ja kontakti sinänsä luontevaa. Välillä potilas unohtaa sanoja, mikä kovasti potilasta tuntuu ärsyttävän. Varsinaisia tunneilmaisuja suuntaan tai toiseen et muutoinjuuri havaitse. Potilas on synkkämielisen oloinen, muttei totaalisen toivoton esim. spontaanisti kertoo odottavansa kesäloman viettoa mökillä.

\begin{itemize}
\tightlist
\item
  \begin{enumerate}
  \def\labelenumi{\alph{enumi})}
  \tightlist
  \item
    Mainitse esitietojen pohjalta neljä erotusdiagnostista tai potilaan vointia mahdollisesti osaltaan selittävää somaattista tekijää ja kerro yhdellä tai kahdella lauseella, miten sulkisit ne pois.
  \end{enumerate}
\item
  \begin{enumerate}
  \def\labelenumi{\alph{enumi})}
  \setcounter{enumi}{1}
  \tightlist
  \item
    Oletetaan, että olet riittävällä varmuudella sulkenut em. syyt pois ja alat kartoittaa tilannetta psykiatrista työdiagnoosia, erotusdiagnostiikkaa ja hoidon suunnittelua varten. Mitä (annettujen esitietojen lisäksi) kysyt potilaalta ja miksi? Listaa kahdeksan mielestäsi tärkeintä asiaa lyhyin perusteluin.
  \end{enumerate}
\item
  \begin{enumerate}
  \def\labelenumi{\alph{enumi})}
  \setcounter{enumi}{2}
  \tightlist
  \item
    Haastattelun jälkeen sinulle on vahvistunut käsitys työdiagnoosista. Miten toteuttaisit alkuvaiheen psykoedukaation potilaalle?
  \end{enumerate}
\end{itemize}

\subsection{Tehtävä 4 (5 pistettä)}\label{tehtuxe4vuxe4-4-5-pistettuxe4}

Äiti saattaa 18-vuotiaan ikätyypillisesti pukeutuneen nuoren naisen yhteispäivystykseen ja kertoo olevansa huolissaan tyttärensä ahdistusoireista. Äiti kertoo, että nuori on ollut hyvin stressaantunut jo viime syksyn kirjoituksista alkaen ja tilanne on koko ajan pahentunut. Nuori on viettänyt jatkuvasti vähemmän aikaa kavereiden kanssa ja viime aikoina ei ole enää poistunut kotoa. Äiti kertoo, että nuori herää vasta vanhempien tullessa töistä kotiin ja vaikuttaa valvovan ja touhuavan jotain öisin. Äiti naurahtaa, että ainahan nuori on sellainen taivaanrannan maalari ollut, mutta nyt on jotenkin aiempaakin poissaolevamman oloinen.
Päätät haastatella nuoren myös ilman äitiä, sillä äidin läsnäollessa nuori antaa käytännössä äidin puhua puolestaan ja lähinnä katselee ympärilleen huoneessa. Kysyt nuorelta alkuun itsetuhoisista ajatuksista, jolloin hän vastaa ''Oon mä kyllä miettinyt elämää ja kuolemaa. Ja linnunrataa. Sitä, et miten se kaikki menee. Ja onks jotain sen takana''. Nuori puhuu hitaasti ja unenomaisesti, pitää välillä pidemmän tauon ja jatkaa ''Ja sitäkin mä oon miettinyt, et miten se kaikki käy järkeen''. Tämän jälkeen nuori vaikuttaa vaipuvan omiin ajatuksiinsa ja hetken kuluttua nousee ylös ja poistuu huoneesta. Palaa huoneeseen kuitenkin aulassa odottaneen äitinsä kanssa.
Keskustelet nuoren ja äidin kanssa hoitovaihtoehdoista. Nuori ei koe olevansa hoidon tarpeessa ja tuo esiin haluavansa takaisin kotiin lukemaan fysiikan kirjoituksiin. Lääkehoitoon liittyvään kysymykseen nuori ei vastaa. Äiti toivoo nuorelle apua ja kertoo, että nuoren saaminen päivystykseenkin oli työn ja tuskan takana.
Jos M1-lähete on mielestäsi tarpeellinen, täytä liitteenä olevaan M1-lähetteeseen kaikki tarpeelliset tiedot siten, että lähete on juridisesti pätevä ja johdonmukainen. Jos M1-lähete ei ole mielestäsi tarpeellinen, kuvaa esseevastauksessa tahdosta riippumattoman hoidon kriteerit ja arvioi kriteerien täyttyminen nuoren kohdalla.

\chapter{2025}\label{section-1}

\chapter{2024}\label{section-2}

\chapter{2023}\label{section-3}

\chapter{2022}\label{section-4}

\chapter{2021}\label{section-5}

\chapter{2020}\label{section-6}

\chapter{2019}\label{section-7}

\chapter{2018}\label{section-8}

\chapter{2017}\label{section-9}

\chapter{2016}\label{section-10}

\chapter{2015}\label{section-11}

\chapter{2014}\label{section-12}

\bibliography{book.bib,packages.bib}

\end{document}
